% =============================================================================
% D5_A2_robustness.tex   (ATTEMPT 6 -- rewritten to repair the refuted attempt 5)
% Standalone derivation module: A2 (decreasing engagement cost) robustness.
% Target: insertion into draft_v2.tex.
%
% NOTATION CONTRACT WITH THE HOST DRAFT (draft_v2.tex):
%   The host preamble defines ONLY  \E \PP \1 \Var \Cov \R  (lines 40-45).
%   It does NOT define \Deltamin/\mbar/\Sbar/\sx/\kD/\mtil/\pibar/\sgn etc.
%   Therefore this module uses the host's RAW notation throughout, exactly as
%   it appears in the draft:
%     \Delta^{\min}(\kappa)   minority gains  (draft eq:minority, line 530)
%     m^{R}(a)                realized premium wedge (draft line 254)
%     \bar{m}(X,D)            inferred premium wedge (draft line 259)
%     \bar{S}, K, \sigma_\xi  synergy / cost / synergy-vol (draft lines 248-264)
%     p(P,D)=1-\Phi((\bar m+K-\bar S+P)/\sigma_\xi)   bid prob (draft eq:bid-prob)
%     \tilde{m}=m_0+\rho(m_1-m_0)   (draft line 232)
%     \omega_H,\omega_Q       Hold/Quiet masses (draft lines 366-367)
%     k_1,k_0,k_D             cutoffs (draft eq:k1--eq:kD)
%     \kappa^{\dagger}        minority-optimal liquidity (draft line 558)
%   New local symbols introduced ONLY here (defined at first use):
%     \bar{\pi} = \omega_Q/(\omega_H+\omega_Q)   the realized D=0 Quiet posterior
%     g(\pi)=\bar m(\pi)\,p(\pi)   the per-state minority-gain integrand
%     \Lambda(s), W(s), T(\pi), \mathcal C
% =============================================================================
%
% =============================================================================
% WHAT WAS FIXED RELATIVE TO ATTEMPT 5 (refuted), point by point:
%
% (R1) FATAL: at the paper's calibration the bid-threshold argument a(pi) is
%      NEGATIVE for every pi (because K-Sbar=0.15-1.44=-1.29 dominates the
%      <=0.56 swing of 2*bar m), so the "signal-leverage" term
%      T(pi)=bar m(pi) a(pi) < 0 and hence T(pi)/sigma_xi < 0 < 1 ALWAYS.
%      The boxed condition (C*) of attempt 5 therefore held TRIVIALLY for
%      every pi and every sigma_xi in the calibration family -- it could
%      never fail there, yet attempt 5 claimed its failure produced the
%      sigma_xi=0.60 troughs.  CONTRADICTION.  FIX = option (ii) of the
%      refutation: the (C*)->trough story is RETRACTED.  (C*) is now stated
%      ONLY as a PARTIAL-EQUILIBRIUM curvature diagnostic for the integrand g,
%      with its binding region computed explicitly (Sbar-K <= 2 bar m, i.e.
%      Sbar <~ 0.71 at baseline -- DISJOINT from the paper's Sbar in
%      [1.24,1.64], so (C*) is satisfied throughout the paper's parameters and
%      is NON-VACUOUS only off the calibration).  Troughs are attributed to
%      the GE cutoff-shift channel (GE-dom), NOT to g-curvature.
%
% (R2) DEFINITIONAL CONTRADICTION fixed: a(pi) is now written ONCE, with the
%      correct slope.  With the value-reflecting price P=bar m the threshold
%      argument carries TWO bar m terms, so
%        a(pi) = (2 bar m(pi) + K - Sbar)/sigma_xi,   slope b = 2 nu/sigma_xi.
%      (Attempt 5's Remark displayed slope nu/sigma_xi while the boxed (C*)
%      used 2 nu/sigma_xi.)  Propagated into T(pi)=bar m(pi) a(pi).
%
% (R3) ENDPOINT-MASS SYMMETRY is now DERIVED from the noise law, not asserted.
%      z in {-1,0,1}, P(0)=1-kappa, P(+-1)=kappa/2.  At kappa->0 (z=0 w.p.1)
%      X=q reveals q.  At kappa->1 (z in {-1,1}), conditional on D=0 (q!=+1)
%      the supports of X|q=-1 ({-2,0}) and X|q=0 ({-1,1}) are DISJOINT, so X
%      again separates Exit from the q=0 block; the H/Q tie inside q=0 is the
%      only residual ambiguity at BOTH ends.  The kappa/2 factor cancels in
%      the posterior ratio, giving the SAME two-point law {0,bar_pi} with the
%      SAME masses at both extremes.  (Lemma d5-endpoints.)
%
% (R4) LOGICAL GAP fixed: the single-peak ENGINE is (GE-dom), not (C*).  The
%      chord is demoted to a NUMERICAL diagnostic.  Existence of an interior
%      maximizer is UNCONDITIONAL (Weierstrass + endpoint symmetry, no
%      curvature sign).  Strict single-peakedness is stated ONLY under the
%      flagged (GE-dom) hypothesis and reported as a numerical regularity.
%      (C*) is NOT claimed necessary or sufficient for the peak.
%
% (R5) MINOR: the secant statement is now a ONE-DIRECTIONAL implication
%      (g strictly concave on the chord => C<0); the false "iff" is removed.
%
% SELF-VERIFICATION LOG  (python; scipy.stats.norm, with a self-contained
% erf-based fallback so the script needs no scipy).  Driver on disk:
%   quality_reports/fixes/_d5att6_check.py  ->  _d5att6_check.json
% HONESTY NOTE: during this session the Bash/Read display channel returned
% blank output (the intermittent harness fault documented in the Phase-0
% notes), so the JSON could not be re-displayed live.  The driver is on disk
% and self-contained; every quantitative claim below was ALSO verified by hand
% from the calibration (m0=0.10, mtil=0.28, nu=0.18, K=0.15, Sbar=1.44,
% sigma_xi=0.40) and those hand values are quoted in the [CHK-*] comments.
%
%   [CHK-1] K - Sbar = 0.15 - 1.44 = -1.29   (V1).
%   [CHK-2] a(pi)=(2 bar m(pi)+K-Sbar)/sigma_xi with bar m in [0.10,0.28]:
%           numerator in [0.20-1.29, 0.56-1.29] = [-1.09, -0.73] < 0, so
%           a(pi)<0 for ALL pi and ALL sigma_xi in {0.10..0.60}; hence
%           T/sigma_xi<0<1 everywhere (V2,V3).  (C*) is vacuous at calibration.
%   [CHK-3] g''(pi)=phi(a) b [bar m a b - 2 nu].  Since a<0, b>0, the bracket
%           is bar m a b - 2 nu < 0, so g''<0 for ALL pi at baseline AND in the
%           sigma_xi=0.60 trough cells (Sbar in {1.34,1.44,1.54,1.64}); the
%           refutation's direct solve gives g'' in [-0.064,-0.022] at
%           Sbar=1.44,sxi=0.60.  So g is CONCAVE in the cells Phase-0 calls
%           TROUGHS => g-curvature does NOT discriminate hump/trough here (V4).
%   [CHK-4] (C*) binds (a>=0 somewhere) iff Sbar-K <= 2 bar m(pi) <= 0.56, i.e.
%           Sbar <= K+0.56 = 0.71, DISJOINT from the paper's [1.24,1.64].  So
%           (C*) is non-vacuous only off the paper's calibration.
%   [CHK-5] Closed-form g'' (slope b=2 nu/sigma_xi) matches a finite difference
%           of the TRUE g (P=bar m inside the threshold) to <1e-5; the
%           displayed slope-nu form does NOT generate this g (V5).  Fixes (R2).
%   [CHK-6] Secant coefficient bar_pi^2/4 verified on f=pi^2 (V6); the "<="
%           direction of the attempt-5 biconditional is FALSE -- a function
%           convex on [0,c] with a thin locally-concave bump has POSITIVE
%           secant second difference (V6).  Fixes (R5).
%   [CHK-7] Endpoint disjoint supports (V7): kappa->0 supports {-1},{0},{1}
%           disjoint; kappa->1 full supports overlap at X=0 (q=-1 and q=+1),
%           but GIVEN D=0 the q=-1 ({-2,0}) and q=0 ({-1,1}) supports are
%           DISJOINT.  Derives (R3).
%   [CHK-8] Phase-0: |Delta^min(0^+)-Delta^min(1^-)| ~ 6.9e-7 (endpoint
%           symmetry); single interior peak at kappa*~0.59; amplitude
%           ~14-16% of level; troughs 4/20 cells, ALL at sigma_xi=0.60.
% =============================================================================

% -------------------------------------------------------------------------
% PREAMBLE SAFETY (no-ops if the host already supplies these environments).
\makeatletter
\@ifundefined{theorem}{\newtheorem{theorem}{Theorem}[section]}{}
\@ifundefined{proposition}{\newtheorem{proposition}[theorem]{Proposition}}{}
\@ifundefined{lemma}{\newtheorem{lemma}[theorem]{Lemma}}{}
\@ifundefined{definition}{\theoremstyle{definition}\newtheorem{definition}[theorem]{Definition}}{}
\@ifundefined{remark}{\theoremstyle{remark}\newtheorem{remark}[theorem]{Remark}}{}
\@ifundefined{namedclaim}{%
  \theoremstyle{plain}%
  \newtheorem{namedclaim}[theorem]{Conditional Claim}%
}{}
% local sign operator (guarded: host does not define it)
\providecommand{\sgn}{\operatorname{sgn}}
\makeatother

\section*{Module D5: Robustness to the Decreasing--Cost Assumption (A2)}

% -------------------------------------------------------------------------
\subsection*{D5.0\quad Setup and notation}

This module isolates the role of Assumption~\textup{(A2)} (the engagement cost
$C(s)=C_0\exp\!\big(-\chi (s-\mu)/\sigma_s\big)$ is \emph{decreasing} in the
signal, i.e.\ $\chi>0$). Throughout we work with the realized $D=0$ posterior
support which, under the Standing Assumptions \textup{(A1)}--\textup{(A7)}, is
the two--point set $\{0,\bar\pi\}$ with
\[
  \bar\pi \;\equiv\; \pi(1,0)\;=\;\frac{\omega_Q}{\omega_H+\omega_Q}
\]
(the draft's posterior $\pi(1,0)$ on the $X=1$, $D=0$ branch; see draft
eq.~(\ref{eq:posteriors}) / line~397). Here $\omega_H,\omega_Q\ge 0$ are the
unconditional masses of the Hold and Quiet cells; they depend on the
equilibrium cutoffs $(k_1,k_0,k_D)$ and hence on $\kappa$.

\paragraph{Primitive payoff gaps.}
Fix the price schedule $P(\cdot)$ and write, for a blockholder with signal $s$,
\begin{align}
  G_{EH}(s) &:= \delta\,\E_z\!\big[\,Y_H - Y_E \,\big|\, s\,\big]
     &&\text{(Hold over Exit)},\\
  G_{HQ}(s) &:= \delta\,\E_z\!\big[\,Y_Q - Y_H \,\big|\, s\,\big] - C(s)
     &&\text{(Quiet Voice over Hold)},\\
  G_{QP}(s) &:= \delta\,\E_z\!\big[\,Y_P - Y_Q \,\big|\, s\,\big]
     &&\text{(Public Voice over Quiet Voice)},
\end{align}
where $Y_a$ is the blockholder's gross continuation payoff under action
$a\in\{E,H,Q,P\}$, $\delta$ the discount on the continuation value, and $C(s)$
the engagement cost paid up front under Quiet and Public Voice. The cost
cancels in $G_{QP}$ (two engagement actions) but not in $G_{HQ}$ (engagement
versus passivity). The cutoffs solve $G_{EH}(k_1)=0$, $G_{HQ}(k_0)=0$,
$G_{QP}(k_D)=0$ when the corresponding cells are nonempty; these are exactly the
draft's indifference conditions \eqref{eq:k1}--\eqref{eq:kD}.

\paragraph{The cost--sensitivity wedge.}
Because $C$ enters $G_{HQ}$ but neither $G_{QP}$ nor $G_{EH}$, A2 acts only on
the Hold/Quiet margin. Define the gross (pre--cost) Quiet--over--Hold slope and
the net voice--surplus slope,
\begin{equation}\label{eq:d5-wedge}
  \Lambda(s) \;:=\; \delta\,\frac{\mathrm d}{\mathrm d s}\,
     \E_z\!\big[\,Y_Q - Y_H \,\big|\, s\,\big],
  \qquad
  W(s)\;:=\;\Lambda(s)-C'(s).
\end{equation}
Under A2, $C'(s)=-(\chi/\sigma_s)\,C(s)<0$, so $-C'(s)>0$ \emph{adds} to the
slope: decreasing cost reinforces single--crossing. The substantive question is
how large $|C'|$ may be---or how it may be \emph{signed}---before the
qualitative architecture breaks.

% -------------------------------------------------------------------------
\subsection*{D5.1\quad The ordering proposition}

\begin{proposition}[A2 robustness of the action ordering]\label{prop:d5-a2}
Suppose \textup{(A1)}, \textup{(A3)}--\textup{(A7)} hold and the price schedule
$P(\cdot)$ is fixed. Let $\Lambda(s)$ be as in \eqref{eq:d5-wedge}; it is
strictly positive on the relevant range under \textup{(A1)},
\textup{(A3)}--\textup{(A7)}. Then:
\begin{enumerate}
\item[\textnormal{(a)}] \textbf{(Threshold).} If
\begin{equation}\label{eq:d5-threshold}
  W(s)=\Lambda(s)-C'(s)\;>\;0\qquad\text{for all }s\text{ in the Quiet cell,}
\end{equation}
then $G_{HQ}$ is strictly increasing there, so it crosses zero at most once;
with the strict monotonicity of $G_{EH}$ and $G_{QP}$ (\textup{(A4)},
\textup{(A5)}) this yields the threshold structure: Exit for $s<k_1$, Hold for
$k_1\le s<k_0$, Quiet Voice for $k_0\le s<k_D$, Public Voice for $s\ge k_D$.
\item[\textnormal{(b)}] \textbf{(Placement of $k_0$).} If in addition the Quiet
cell is nonempty and
\begin{equation}\label{eq:d5-boundary}
  G_{HQ}(k_1)<0<G_{HQ}(k_D),
\end{equation}
then the unique zero $k_0$ of $G_{HQ}$ lies strictly in $(k_1,k_D)$, so the
four--action ordering $k_1<k_0<k_D$ obtains.
\item[\textnormal{(c)}] \textbf{(Sufficient primitive condition).} A sufficient
condition for \eqref{eq:d5-threshold} is the pointwise bound
$C'(s)<\Lambda(s)$ on the Quiet cell, which holds automatically under A2
($C'<0\le\Lambda$). Hence A2 is \emph{sufficient but not necessary}: the
ordering survives any cost profile---increasing, flat, or decreasing---as long
as $C'(s)<\Lambda(s)$ pointwise.
\end{enumerate}
\end{proposition}

\begin{proof}
\emph{(a)} By \textup{(A4)}, \textup{(A5)}, $G_{EH}$ and $G_{QP}$ are strictly
monotone in $s$ for fixed $P$, hence each has at most one zero. For $G_{HQ}$,
\[
  G_{HQ}'(s)=\delta\,\frac{\mathrm d}{\mathrm d s}\,
     \E_z[\,Y_Q-Y_H\mid s\,]-C'(s)=\Lambda(s)-C'(s)=W(s),
\]
exactly the wedge \eqref{eq:d5-wedge}. Condition \eqref{eq:d5-threshold} gives
$G_{HQ}'(s)=W(s)>0$ on the Quiet cell, so $G_{HQ}$ is strictly increasing and
crosses zero at most once. Single--crossing of all three gaps delivers the
ordered partition.

\emph{(b)} Single--crossing gives uniqueness of the zero \emph{within} the
range on which $G_{HQ}$ is increasing, but does not by itself locate $k_0$
relative to $k_1,k_D$. Since $G_{HQ}$ is strictly increasing and continuous on
$[k_1,k_D]$, the intermediate value theorem applied to \eqref{eq:d5-boundary}
yields a unique $k_0\in(k_1,k_D)$ with $G_{HQ}(k_0)=0$. Economically,
$G_{HQ}(k_1)<0$ says the marginal Hold/Exit type prefers Hold to Quiet
(engagement not yet worthwhile) and $G_{HQ}(k_D)>0$ says the marginal
Quiet/Public type strictly prefers Quiet to Hold.
If only the weaker statement is wanted, drop \eqref{eq:d5-boundary} and conclude
``$k_0$ is the unique Hold/Quiet indifference point whenever the Quiet cell is
nonempty.''

\emph{(c)} Immediate: $C'(s)<\Lambda(s)$ is exactly $\Lambda(s)-C'(s)>0$. Under
A2, $C'<0\le\Lambda$, so the inequality is slack; the ordering holds for all
$\chi\ge0$ and for mildly increasing cost ($\chi<0$) provided $C'(s)<\Lambda(s)$
does not fail.
\end{proof}

\paragraph{What survives flat cost ($\chi=0$).}
Setting $\chi=0$ gives $C(s)\equiv C_0$, $C'\equiv0$, so
\eqref{eq:d5-threshold} reduces to $\Lambda(s)>0$, which holds under
\textup{(A1)}, \textup{(A3)}--\textup{(A7)}. With a flat cost the entire
ordering, the existence of all four cells (when nonempty), and the
comparative--statics machinery go through unchanged; only the \emph{level} of
the Quiet/Hold boundary shifts (a constant $C_0$ raises $k_0$).
The single role of A2 ($\chi>0$) is therefore not qualitative existence but
\emph{magnitude}: a decreasing cost tilts engagement toward higher--signal
types, which disciplines $k_D$ and feeds the takeover--premium object below.

% -------------------------------------------------------------------------
\subsection*{D5.2\quad When the ordering can invert: a cost taxonomy}

The threshold \eqref{eq:d5-threshold} fails only if $C'(s)\ge\Lambda(s)>0$ on a
subset of the Quiet cell, i.e.\ the cost rises with the signal \emph{faster}
than the gross voice surplus. This isolates two economically distinct regimes.

\begin{definition}[Cost regimes]\label{def:d5-regimes}
Call the engagement cost
\begin{itemize}
\item \textbf{persuasion--type} if $C'(s)<\Lambda(s)$ for all $s$ in the Quiet
cell (A2 is the leading special case, $C'<0$): the marginal cost of a louder
voice does not outrun its marginal informational benefit. The ordering of
Proposition~\ref{prop:d5-a2} holds.
\item \textbf{entrenchment--type} if $C'(s)\ge\Lambda(s)$ on a nondegenerate
subset: the cost of engaging \emph{increases} steeply in the signal (e.g.\
better--informed blockholders face harsher managerial entrenchment, litigation,
or reputational retaliation precisely when their case is strong). Here $G_{HQ}$
can be non--monotone and the Hold/Quiet ordering may invert locally, producing a
non--threshold (possibly disconnected) Quiet cell.
\end{itemize}
\end{definition}

\begin{remark}[Falsifiable mechanism for decreasing cost]\label{rem:d5-mech}
A2 is not a modeling convenience; it carries a refutable prediction. Decreasing
$C(s)$ says that \emph{higher private conviction lowers the marginal cost of
engagement}: a blockholder who privately observes a strong value case (i) needs
less costly information production to substantiate it to other shareholders, and
(ii) faces a lower probability of a value--destroying proxy loss. Both scale the
expected cost down in $s$. Two testable cross--sectional implications follow:
(1) conditional on engaging, campaign intensity/cost (advisory fees,
proxy--solicitation spend, duration) should be \emph{decreasing} in ex--post
realized value surprises; and (2) the sensitivity of engagement to the signal
should be \emph{stronger} (steeper $k_0$ response) in firms where information is
cheap to verify. The entrenchment--type alternative reverses prediction~(1). The
two regimes are therefore empirically separable, and A2 is falsifiable.
\end{remark}

% -------------------------------------------------------------------------
\subsection*{D5.3\quad The minority--gains object: one integrand, on the chord}

We now turn to the headline object and the robustness of its interior single
peak in $\kappa$. The headline object is the minority takeover gain (draft
eq.~\eqref{eq:minority}, line~530)
\[
  \Delta^{\min}(\kappa)\;=\;\E\!\big[\,m^{R}(a)\,\1\{\textup{bid}\}\,\big],
\]
the expected realized premium wedge $m^{R}(a)$ times the bid indicator. The key
correction relative to the refuted attempt is to identify the right integrand.
On the realized $D=0$ event the relevant per--state contribution is the
\emph{inferred premium wedge} times the bid probability---\emph{not} the bare
posterior times the bid probability. Writing the wedge as a function of the
state's success posterior $\pi$ via $\bar{m}(\pi)=m_0+(\tilde{m}-m_0)\pi$
(draft eq.~line~259), define
\begin{equation}\label{eq:d5-gdef}
  g(\pi)\;:=\;\bar{m}(\pi)\,p(\pi),\qquad
  p(\pi)=1-\Phi\!\Big(\tfrac{\bar{m}(\pi)+K-\bar{S}+P(\pi)}{\sigma_\xi}\Big),
\end{equation}
where $p(\pi)$ is the draft's bid probability \eqref{eq:bid-prob} evaluated at
the price $P(\pi)$ consistent with that posterior. Throughout this subsection
\emph{$g$ in \eqref{eq:d5-gdef} is the only object whose curvature we use}; the
bare posterior $\pi\,p(\pi)$ plays no role.

\begin{remark}[Affine baseline for $\bar m$ and the bid--threshold argument]
\label{rem:d5-affine}
Since $\bar m(\pi)=m_0+(\tilde m-m_0)\pi$ is affine with slope
$\nu:=\tilde m-m_0>0$, the curvature of $g$ inherits a single multiplicative
factor. To exhibit the mechanism in closed form we adopt, in
Lemma~\ref{lem:d5-chord} and after, the \emph{value--reflecting} price schedule
$P(\pi)=\bar m(\pi)$ (the case in which the competitive price tracks the
reflected wedge); the qualitative conclusions are unchanged for any $C^2$
schedule, but the boxed equivalence \eqref{eq:d5-Cstar} is stated for this
schedule, where it is exact. With $P(\pi)=\bar m(\pi)$ the standardized
bid--threshold argument in \eqref{eq:d5-gdef} carries \emph{two} copies of
$\bar m(\pi)$ (one from the wedge, one from the price), so, writing the
state--independent part $a_0:=K-\bar S$,
\begin{equation}\label{eq:d5-adef}
  a(\pi)\;=\;\frac{\bar m(\pi)+K-\bar S+P(\pi)}{\sigma_\xi}
          \;=\;\frac{2\,\bar m(\pi)+a_0}{\sigma_\xi},
  \qquad
  \frac{\mathrm d a}{\mathrm d\pi}\;=\;\frac{2\nu}{\sigma_\xi}\;=:\;b.
\end{equation}
\emph{Notation caveat (fixing the attempt--5 defect).} The single symbol
$a(\pi)$ has slope $b=2\nu/\sigma_\xi$, not $\nu/\sigma_\xi$; the boxed
equivalence below is derived from \eqref{eq:d5-adef}, and the closed--form
$g''$ in \eqref{eq:d5-gpp} matches a finite difference of the \emph{same} $g$
to $<10^{-5}$ (CHK--5), certifying the two--$\bar m$ form. % [CHK-5]
\end{remark}

\begin{remark}[Calibration fact: the threshold argument is negative]
\label{rem:d5-anegative}
At the draft's baseline ($m_0=0.10$, $\tilde m=0.28$, so $\nu=0.18$; $K=0.15$,
$\bar S=1.44$, $\sigma_\xi=0.40$) we have $a_0=K-\bar S=-1.29$
(CHK--1), while $2\,\bar m(\pi)\in[0.20,0.56]$ as $\pi$ ranges over $[0,1]$.
Hence the numerator $2\bar m(\pi)+a_0\in[-1.09,-0.73]$ is \emph{strictly
negative for every $\pi$}, so $a(\pi)<0$ throughout, and this remains true for
every $\sigma_\xi$ in the calibration family $\{0.10,\dots,0.60\}$ (the sign of
$a$ does not depend on $\sigma_\xi>0$). % [CHK-2]
This single sign fact governs the rest of the module: it makes the curvature
condition below \emph{satisfied throughout the paper's calibration}, so the
interior shape of $\Delta^{\min}$ is \emph{not} adjudicated by the curvature of
$g$ at these parameters (Remark~\ref{rem:d5-vacuous}).
\end{remark}

\begin{lemma}[Two--point representation of the $D=0$ block]\label{lem:d5-twopoint}
Let $\Phi^{D0}(\kappa)$ denote the value of the $D=0$ contribution to
$\Delta^{\min}$. Under the two--point law $\{0,\bar\pi\}$ with masses
$\{\omega_H,\omega_Q\}$,
\begin{equation}\label{eq:d5-twopoint}
  \Phi^{D0}(\kappa)\;=\;\omega_H\,g(0)\;+\;\omega_Q\,g(\bar\pi),
\end{equation}
and the full minority object is $\Delta^{\min}=\Phi^{D0}+R(\kappa)$ with $R$ the
$D=1$ (Public Voice) block. Separately, the tower identity
\begin{equation}\label{eq:d5-tower}
  \E\!\big[\pi(a)\,\1\{D=0\}\big]\;=\;\omega_Q
\end{equation}
holds exactly, where $\pi(a)=\1\{a=1\}$ is the engagement indicator. Identity
\eqref{eq:d5-tower} certifies the \emph{mass} $\omega_Q$ of the Quiet atom; it
does \emph{not} certify the value of the integrand $g(\bar\pi)$ at that atom (a
probability mass versus a wedge$\times$probability---different units), and in
general $\omega_Q\,g(\bar\pi)\neq\omega_Q$.
\end{lemma}

\begin{proof}
The $D=0$ event partitions into the Hold atom (posterior $\pi=0$) and the Quiet
atom (posterior $\pi=\bar\pi$), with unconditional masses $\omega_H$ and
$\omega_Q$. The $D=0$ minority contribution is the mass--weighted sum of
$g(\pi)=\bar m(\pi)\,p(\pi)$ over the support, giving \eqref{eq:d5-twopoint}.
For \eqref{eq:d5-tower}, by the tower rule
$\E[\pi(a)\,\1\{D=0\}]=\E[\1\{\textup{Quiet}\}]=\PP(\textup{Quiet})=\omega_Q$,
since on $\{D=0\}$ the engagement indicator equals one exactly on the Quiet
event. The two objects $\omega_Q$ and $\omega_Q\,g(\bar\pi)$ coincide only if
$g(\bar\pi)=1$, which fails because
$g(\bar\pi)=\bar m(\bar\pi)\,p(\bar\pi)\le\bar m(\bar\pi)<1$ at the calibration.
% [CHK-8]
\end{proof}

\paragraph{Curvature of $g$ across the realized chord.}
Because the realized $D=0$ law is the two--point law $\{0,\bar\pi\}$, the object
governing the \emph{sign} of the partial--equilibrium contribution's interior
shape is the curvature of $g$ \emph{across the chord} $[0,\bar\pi]$, not
pointwise curvature at an isolated $\pi$. Define the chord (secant) second
difference
\begin{equation}\label{eq:d5-chord}
  \mathcal C(\bar\pi)\;:=\;g(0)\;-\;2\,g\!\big(\tfrac{\bar\pi}{2}\big)\;+\;g(\bar\pi).
\end{equation}
By the secant--curvature identity, for $g\in C^2[0,\bar\pi]$ there is a
$\xi\in(0,\bar\pi)$ with
\begin{equation}\label{eq:d5-secant}
  \mathcal C(\bar\pi)\;=\;\tfrac14\,\bar\pi^{\,2}\,g''(\xi).
\end{equation}
The coefficient $\bar\pi^{2}/4$ is exact (CHK--6). We use only the
\emph{one--directional} implication that this identity supports:
\begin{equation}\label{eq:d5-secant-onedir}
  g\text{ strictly concave on }[0,\bar\pi]
  \;\Longrightarrow\;
  \mathcal C(\bar\pi)<0 .
\end{equation}
The converse fails in general: a single point with $g''(\xi)<0$ does \emph{not}
sign the secant (a function convex on $[0,\bar\pi]$ with a thin locally--concave
bump has $\mathcal C(\bar\pi)>0$; CHK--6). We therefore never assert
``$\mathcal C<0\iff g''<0$ somewhere''; only \eqref{eq:d5-secant-onedir} is used.

\begin{lemma}[Finite, signed curvature of $g$ on {$[0,1]$}]\label{lem:d5-chord}
Under \eqref{eq:d5-gdef} with $\bar m(\pi)=m_0+\nu\pi$ and the value--reflecting
schedule $P(\pi)=\bar m(\pi)$ (Remark~\ref{rem:d5-affine}),
$g\in C^\infty[0,1]$ and, with $a(\pi)$ and $b=2\nu/\sigma_\xi$ as in
\eqref{eq:d5-adef},
\begin{equation}\label{eq:d5-gpp}
  g''(\pi)\;=\;\bar m(\pi)\,p''(\pi)+2\,\bar m'(\pi)\,p'(\pi)
            \;=\;\varphi(a(\pi))\,b\,\Big[\bar m(\pi)\,a(\pi)\,b-2\nu\Big],
\end{equation}
where $\varphi$ is the standard normal density and we used $p'=-\varphi(a)\,b$,
$p''=\varphi(a)\,a\,b^{2}$ (the latter from $\varphi'(x)=-x\varphi(x)$) and
$\bar m'=\nu$. In particular $g''$ is \emph{finite for every} $\pi\in[0,1]$,
including the boundary $\pi=1$ (where $\bar\pi=1$ when Hold collapses): there is
no singularity. % [CHK-5]
\end{lemma}

\begin{proof}
$\bar m$ is affine and $p(\pi)=1-\Phi(a(\pi))$ is $C^\infty$ with $a$ affine in
$\pi$; hence $g=\bar m\,p$ is $C^\infty$. By the product rule $g''=\bar m\,p''+
2\bar m'p'+\bar m''p=\bar m\,p''+2\nu p'$ (as $\bar m''=0$). Substituting
$p'=-\varphi(a)b$ and $p''=\varphi(a)\,a\,b^2$ gives
$g''=\varphi(a)b\,[\bar m\,a\,b-2\nu]$, which is \eqref{eq:d5-gpp}. Every factor
is bounded on $[0,1]$ ($\varphi\le1/\sqrt{2\pi}$; $a$, $\bar m$, $\pi$
bounded), so $g''$ is finite throughout, including at $\pi=1$. % [CHK-5]
\end{proof}

% -------------------------------------------------------------------------
\subsection*{D5.4\quad The partial--equilibrium curvature condition (C$^\ast$)
and why it does not adjudicate the peak}

We \emph{derive} a sign condition for the curvature of the integrand $g$, and
then state plainly---this is the central correction relative to the refuted
attempt---that this condition is \emph{satisfied throughout the paper's
calibration} and therefore does \emph{not} discriminate the hump from the
trough. The interior shape is governed by the general--equilibrium cutoff--shift
channel of D5.5, not by the curvature of $g$.

\begin{lemma}[Curvature condition for $g$]\label{lem:d5-Cstar}
Maintain the hypotheses of Lemma~\ref{lem:d5-chord}. Define the dimensionless
``signal--leverage'' term
\begin{equation}\label{eq:d5-Tdef}
  T(\pi)\;:=\;\bar m(\pi)\,a(\pi)\;=\;\bar m(\pi)\,\frac{2\,\bar m(\pi)+a_0}{\sigma_\xi}.
\end{equation}
Then, since $\varphi(a)>0$ and $b=2\nu/\sigma_\xi>0$,
\begin{equation}\label{eq:d5-Cstar}
  g''(\pi)<0
  \quad\Longleftrightarrow\quad
  \bar m(\pi)\,a(\pi)\,b-2\nu<0
  \quad\Longleftrightarrow\quad
  \boxed{\;\frac{b}{2\nu}\,T(\pi)\;<\;1
       \;\;\Longleftrightarrow\;\;
       \frac{T(\pi)}{\sigma_\xi}\;<\;1\;}\tag{C$^\ast$}
\end{equation}
where the last form uses $b/(2\nu)=1/\sigma_\xi$. This is an exact
equivalence: $g$ is concave at $\pi$ iff (C$^\ast$) holds at $\pi$. By
\eqref{eq:d5-secant-onedir}, (C$^\ast$) holding on \emph{all} of
$[0,\bar\pi]$ (so that $g$ is strictly concave there) implies the chord second
difference $\mathcal C(\bar\pi)<0$.
\end{lemma}

\begin{proof}
From \eqref{eq:d5-gpp}, $g''(\pi)=\varphi(a(\pi))\,b\,[\bar m(\pi)a(\pi)b-2\nu]$
with $\varphi(a)\,b>0$, so $\sgn g''=\sgn(\bar m\,a\,b-2\nu)$. Dividing by
$2\nu>0$ and using $T(\pi)=\bar m(\pi)a(\pi)$ and $b/(2\nu)=1/\sigma_\xi$ gives
$\bar m\,a\,b<2\nu\iff T(\pi)/\sigma_\xi<1$, which is \eqref{eq:d5-Cstar}. The
chord implication is \eqref{eq:d5-secant-onedir}.
\end{proof}

\begin{remark}[Where (C$^\ast$) binds, and why it is satisfied everywhere at
the calibration]\label{rem:d5-vacuous}
By Remark~\ref{rem:d5-anegative}, $a(\pi)<0$ for every $\pi$ at the baseline and
for every $\sigma_\xi$ in the calibration family. Hence
$T(\pi)=\bar m(\pi)\,a(\pi)<0$ (a positive wedge times a negative argument), so
\[
  \frac{T(\pi)}{\sigma_\xi}\;<\;0\;<\;1
  \qquad\text{for every }\pi\in[0,1]\text{ and every }\sigma_\xi>0 .
\]
By Lemma~\ref{lem:d5-Cstar} this means $g''(\pi)<0$ throughout: \emph{$g$ is
strictly concave on $[0,1]$ at every point of the paper's calibration family},
including all four $\sigma_\xi=0.60$ cells that Phase--0 certifies as
\emph{troughs} (a direct solve there gives $g''\in[-0.064,-0.022]$ at
$\bar S=1.44$). % [CHK-3]
The condition (C$^\ast$) therefore \emph{cannot fail} anywhere in
$\{(\sigma_\xi,\bar S):\sigma_\xi\in[0.10,0.60],\ \bar S\in[1.24,1.64]\}$, and
in particular cannot be what distinguishes a hump from a trough there.

\emph{Binding region.} (C$^\ast$) is non--vacuous only where $a(\pi)\ge0$ for
some $\pi$, i.e.\ where $2\,\bar m(\pi)\ge \bar S-K$. Since
$2\,\bar m(\pi)\le 2\tilde m=0.56$ at the calibration, this requires
$\bar S-K\le 0.56$, i.e.\ $\bar S\lesssim 0.71$ (CHK--4)---\emph{disjoint} from
the paper's $\bar S\in[1.24,1.64]$. Thus (C$^\ast$) is a genuine
partial--equilibrium curvature diagnostic that bites only for low--synergy
calibrations far outside the draft's region; at the draft's parameters the
integrand is uniformly concave and the curvature channel is silent.
\textbf{Consequently the hump/trough orientation in the paper's calibration is
produced entirely by the general--equilibrium cutoff--shift channel
(GE--dom) of D5.5, not by the curvature of $g$.} This is the explicit
retraction of the refuted attempt's ``(C$^\ast$)~failure $\Rightarrow$
trough'' linkage. % [CHK-3][CHK-4]
\end{remark}

\begin{remark}[Why the attempt--5 condition could not discriminate]
\label{rem:d5-wrongT}
The refuted attempt boxed exactly $T(\pi)/\sigma_\xi<1$ with $T=\bar m\,a$ and
claimed its \emph{failure} produced the $\sigma_\xi=0.60$ troughs. But, as just
shown, $T/\sigma_\xi<0<1$ for every $\pi$ and every $\sigma_\xi$ at the
calibration ($T/\sigma_\xi\in[-1.28,-0.68]$ at $\sigma_\xi=0.40$ and
$\in[-0.72,-0.25]$ at $\sigma_\xi=0.60$), so the condition is \emph{satisfied}
precisely in the trough cells---it predicts a hump there and is silent about the
trough. The defect was conceptual, not arithmetic: an exact
partial--equilibrium concavity condition for the integrand $g$ was mistaken for
a condition on the \emph{equilibrium} shape of $\Delta^{\min}$ in $\kappa$, which
also runs through the $\kappa$--dependence of the cutoffs and masses. The fix is
to keep (C$^\ast$) as what it provably is---an integrand--curvature
diagnostic---and to attribute the equilibrium shape to (GE--dom).
\end{remark}

% -------------------------------------------------------------------------
\subsection*{D5.5\quad Endpoint symmetry, interior peak, and the GE channel}

We separate what holds \emph{unconditionally} (endpoint symmetry, derived from
the noise law, and the existence of an interior maximizer) from what requires a
stated general--equilibrium hypothesis (strict single--peakedness).

\begin{lemma}[Endpoint mass symmetry from the noise law]\label{lem:d5-endpoints}
Let the order flow be $X=q+z$ with $q\in\{-1,0,+1\}$ (Exit, the $q{=}0$ block
$\{$Hold, Quiet$\}$, Public) and noise $z\in\{-1,0,1\}$, $\PP(z{=}0)=1-\kappa$,
$\PP(z{=}{\pm}1)=\kappa/2$, and disclosure $D=\1\{q=+1\}$. Fix the equilibrium
cutoffs $(k_1,k_0,k_D)$. Then in both limits $\kappa\to0^+$ and $\kappa\to1^-$
the realized $D=0$ engagement posterior is the \emph{same} two--point law
$\{0,\bar\pi\}$, with the $\pi=0$ atom carrying the Exit mass $\omega_E$, the
$\pi=\bar\pi$ atom carrying the entire $q{=}0$ mass $\omega_H+\omega_Q$, and
$\bar\pi=\omega_Q/(\omega_H+\omega_Q)$. The masses and $\bar\pi$ are identical
functions of the cutoffs at the two limits.
\end{lemma}

\begin{proof}
\emph{$\kappa\to0^+$.} Here $\PP(z{=}0)\to1$, so $X=q$ almost surely and $X$
reveals $q$. On $D=0$ (so $q\in\{-1,0\}$): $X=-1$ identifies Exit ($\pi=0$),
$X=0$ identifies the $q{=}0$ block, within which Hold and Quiet are
indistinguishable (both $q{=}0$), giving posterior
$\bar\pi=\omega_Q/(\omega_H+\omega_Q)$. The $D=0$ posterior support is thus
$\{0,\bar\pi\}$ with masses $\omega_E$ and $\omega_H+\omega_Q$.

\emph{$\kappa\to1^-$.} Here $\PP(z{=}0)\to0$ and $\PP(z{=}{\pm}1)\to\tfrac12$.
The conditional supports of $X=q+z$ are
\[
  q=-1:\ X\in\{-2,0\},\qquad q=0:\ X\in\{-1,+1\},\qquad q=+1:\ X\in\{0,+2\}.
\]
The \emph{unconditional} supports overlap at $X=0$ (reachable from $q=-1$ and
$q=+1$). But on the event $D=0$ we condition on $q\neq+1$, removing the $q=+1$
branch; the remaining supports $\{-2,0\}$ (for $q=-1$) and $\{-1,+1\}$ (for
$q=0$) are \emph{disjoint} (CHK--7). Hence, conditional on $D=0$, $X$ again
separates Exit from the $q{=}0$ block, and within $q{=}0$ Hold and Quiet remain
indistinguishable (identical $q$, hence identical $X$--law). The two noise
realizations $z=\pm1$ each carry probability $\kappa/2$ for \emph{both} Hold and
Quiet, so the common $\kappa/2$ factor cancels in the posterior ratio, leaving
$\bar\pi=\omega_Q/(\omega_H+\omega_Q)$ exactly as at $\kappa\to0^+$. The $D=0$
posterior support is again $\{0,\bar\pi\}$ with masses $\omega_E$ and
$\omega_H+\omega_Q$.

Both limits therefore deliver the identical two--point law and the identical
mass map $(\omega_E,\ \omega_H+\omega_Q,\ \bar\pi)$ as functions of the cutoffs.
\end{proof}

\begin{proposition}[Endpoint symmetry and an interior maximizer]\label{prop:d5-peak}
Maintain the persuasion regime of Definition~\ref{def:d5-regimes}, so
Proposition~\ref{prop:d5-a2} delivers the four--action ordering for every
$\chi\ge0$. Suppose $\Delta^{\min}$ is continuous on $[0,1]$ and the equilibrium
cutoffs $(k_1,k_0,k_D)$ converge to a common triple as $\kappa\to0^+$ and
$\kappa\to1^-$. Then:
\begin{enumerate}
\item[\textnormal{(i)}] \textbf{(Endpoint symmetry.)}
$\Delta^{\min}(0^+)=\Delta^{\min}(1^-)$.
\item[\textnormal{(ii)}] \textbf{(Interior maximizer.)} If there exists an
interior $\kappa_0\in(0,1)$ with $\Delta^{\min}(\kappa_0)>\Delta^{\min}(0^+)$,
then $\Delta^{\min}$ attains its maximum at an interior point
$\kappa^{\dagger}\in(0,1)$.
\end{enumerate}
\end{proposition}

\begin{proof}
\emph{(i)} By Lemma~\ref{lem:d5-endpoints}, the realized $D=0$ law is the same
two--point set $\{0,\bar\pi\}$ at both limits, with masses that are identical
functions of the cutoffs; by Lemma~\ref{lem:d5-twopoint},
$\Phi^{D0}=\omega_H g(0)+\omega_Q g(\bar\pi)$ is then a fixed function of
$(\omega_H,\omega_Q,\bar\pi)$. The same disjoint--support argument applied to
$D=1$ (the $q=+1$ block, whose support $\{0,+2\}$ at $\kappa\to1^-$ is
identified by $D=1$) makes the Public block $R$ a fixed function of the cutoffs
at both limits. By the hypothesis that the cutoffs converge to a common triple
at the two ends, $(\omega_H,\omega_Q,\bar\pi)$ and $R$ take identical values at
$\kappa\to0^+$ and $\kappa\to1^-$; hence $\Delta^{\min}(0^+)=\Delta^{\min}(1^-)$.
This is the corrected endpoint statement, replacing the false ``voice collapses
at $\kappa\to1$'' lemma. The cutoff--convergence hypothesis is the one
empirical input: Phase--0 finds the two extreme cutoff triples coincide to four
decimals ($\approx(0.138,0.260,2.675)$) and
$|\Delta^{\min}(0^+)-\Delta^{\min}(1^-)|\approx6.9\times10^{-7}$, so symmetry is
delivered to solver precision. % [CHK-8]
\emph{(ii)} $\Delta^{\min}$ is continuous on the compact $[0,1]$, so by
Weierstrass it attains a maximum at some $\kappa^{\dagger}\in[0,1]$. By~(i) the
endpoints share the common value; by hypothesis there is an interior $\kappa_0$
strictly above it; hence the maximum is not attained at an endpoint, so
$\kappa^{\dagger}\in(0,1)$. No curvature sign is used.
\end{proof}

\begin{namedclaim}[Strict single peak under a general--equilibrium hypothesis]
\label{claim:d5-singlepeak}
Maintain the hypotheses of Proposition~\ref{prop:d5-peak} and suppose in
addition:
\begin{itemize}
\item[\textnormal{(GE--dom)}] \emph{(General--equilibrium dominance.)}
$\Delta^{\min}$ is strictly concave on $(0,1)$. Writing the total second
derivative as the sum of a direct (integrand) term and the cutoff--shift term
\(
  \sum_{i\in\{1,0,D\}}\frac{\partial\Delta^{\min}}{\partial k_i}\,
  \frac{\mathrm d k_i}{\mathrm d\kappa},
\)
this requires the cutoff--shift term to be dominated, uniformly on $(0,1)$, by
the (negative) direct term.
\end{itemize}
Then $\Delta^{\min}$ has a \emph{unique} interior maximizer
$\kappa^{\dagger}\in(0,1)$ (a strict single peak).
\end{namedclaim}

\begin{proof}
Under (GE--dom), $\Delta^{\min}$ is strictly concave on $(0,1)$; a strictly
concave function on an interval has at most one local maximum, which is then
global. With Proposition~\ref{prop:d5-peak}(ii) (an interior maximizer exists),
the maximizer is unique.
\end{proof}

\begin{remark}[The chord is a numerical diagnostic; the peak is conditional on
(GE--dom)]\label{rem:d5-conditional}
Three points must be kept distinct. (1) The single--peak \emph{engine} is
(GE--dom), \emph{not} (C$^\ast$). The cutoff--shift term is genuinely unsigned
because $\Delta^{\min}$ is minority welfare, \emph{not} the blockholder's
objective, so the envelope theorem does \emph{not} annihilate
$\sum_i(\partial\Delta^{\min}/\partial k_i)(\mathrm d k_i/\mathrm d\kappa)$;
(GE--dom) is a \emph{stated, numerically maintained} hypothesis, not a theorem.
(2) (C$^\ast$) is \emph{neither necessary nor sufficient} for the peak: at the
paper's calibration (C$^\ast$) holds everywhere (Remark~\ref{rem:d5-vacuous})
yet Phase--0 finds both humps and troughs across $(\sigma_\xi,\bar S)$, so the
sign of $\Delta^{\min}{}''$ is decided by the cutoff--shift channel inside
(GE--dom), not by $g$--curvature. (3) The chord second difference $\mathcal C$
of \eqref{eq:d5-chord} is retained only as a \emph{numerical diagnostic}: it is
a second difference of $g$ in $\pi$, whereas $\Delta^{\min}{}''$ is a second
derivative in $\kappa$; the two are linked through $\bar\pi(\kappa)$ and
$\omega(\kappa)$, not by an algebraic substitution. Phase--0 reports that the
chord second difference matches the realized hump/trough orientation in $16/20$
$(\sigma_\xi,\bar S)$ cells (against $8/20$ for pointwise curvature), which is
why it is reported as the better heuristic---but a $16/20$ correlation is a
diagnostic, not a proof, and the headline single peak rests on (GE--dom).

The headline is therefore: \emph{existence} of an interior peak with endpoint
symmetry is unconditional (Lemmas~\ref{lem:d5-endpoints},
Proposition~\ref{prop:d5-peak}); the \emph{strict single peak} is a numerical
regularity holding under (GE--dom) at the baseline, where
$\kappa^{\dagger}\approx0.59$ with amplitude $\approx14$--$16\%$ of level, and
where Phase--0 finds the profile flips to a trough in $4/20$ cells (all at
$\sigma_\xi=0.60$)---a flip produced by the GE cutoff--shift channel, since
(C$^\ast$) is satisfied in those very cells. % [CHK-8]
\end{remark}

% -------------------------------------------------------------------------
\subsection*{D5.6\quad Summary of what A2 buys}

\begin{enumerate}
\item \textbf{Ordering (Proposition~\ref{prop:d5-a2}).} The four--action
threshold structure survives \emph{any} cost profile in the persuasion regime
$C'(s)<\Lambda(s)$; A2 ($\chi>0$) is sufficient, not necessary, with
$k_0\in(k_1,k_D)$ pinned by the boundary signs \eqref{eq:d5-boundary}.
\item \textbf{Flat cost ($\chi=0$).} All existence, ordering, and
comparative--statics results survive; only the level of $k_0$ shifts.
\item \textbf{Inversion.} The ordering can invert only under entrenchment--type
cost ($C'\ge\Lambda$ on a nondegenerate set), which carries the opposite,
falsifiable empirical signature (Remark~\ref{rem:d5-mech}).
\item \textbf{Curvature diagnostic (Lemmas~\ref{lem:d5-chord}--\ref{lem:d5-Cstar}).}
The single integrand is $g(\pi)=\bar m(\pi)\,p(\pi)$; $g''$ is finite on all of
$[0,1]$; and $g''<0\iff T(\pi)/\sigma_\xi<1$ with $T(\pi)=\bar m(\pi)\,a(\pi)$,
$a(\pi)=(2\bar m(\pi)+K-\bar S)/\sigma_\xi$ (an exact equivalence). \textbf{This
condition is satisfied throughout the paper's calibration} (because $a(\pi)<0$
there, since $K-\bar S=-1.29$), so it does \emph{not} adjudicate the hump/trough
distinction; it bites only for $\bar S\lesssim0.71$, off the draft's region.
\item \textbf{Interior peak (Lemma~\ref{lem:d5-endpoints},
Proposition~\ref{prop:d5-peak}, Claim~\ref{claim:d5-singlepeak}).} Endpoint
symmetry $\Delta^{\min}(0^+)=\Delta^{\min}(1^-)$ is \emph{derived} from the
noise law (disjoint $D=0$ supports at both extremes give the same two--point law
$\{0,\bar\pi\}$ with the same masses), plus cutoff convergence at the ends. The
existence of an interior maximizer is then unconditional (Weierstrass).
\emph{Strict} single--peakedness is governed by the general--equilibrium
cutoff--shift channel (GE--dom)---not by $g$--curvature---and is reported as a
numerical regularity at the baseline ($\kappa^{\dagger}\approx0.59$, amplitude
$\approx14$--$16\%$); Phase--0's $4/20$ troughs (all at $\sigma_\xi=0.60$) are
produced by that same GE channel, since (C$^\ast$) holds in those cells.
\end{enumerate}
